% Bagian: first-order-logic
% Bab: model-theory
% Subbagian: theory-of-m

\documentclass[../../../include/open-logic-section]{subfiles}

\begin{document}

\olfileid[id]{mod}{bas}{thm}

\section{Teori suatu \printtoken{S}{structure}}

Setiap !!{structure}~$\Struct{M}$ membuat sebagian !!{sentence}s benar dan
sebagian lainnya salah. Himpunan semua !!{sentence}s yang dibuatnya benar
disebut \emph{teori} struktur tersebut. Himpunan itu memang merupakan suatu
teori, karena apa pun yang diakibatkannya harus benar dalam semua modelnya,
termasuk~$\Struct{M}$.

\begin{defn}
  Diberikan !!a{structure}~$\Struct M$, \emph{teori} dari $\Struct{M}$
  ialah himpunan $\Theory{M}$ dari !!{sentence}s yang benar dalam
  $\Struct{M}$, yakni $\Theory{M} = \Setabs{!A}{\Sat{M}{!A}}$.
\end{defn}

Kita juga menggunakan istilah ``teori'' secara informal untuk mengacu pada
himpunan-himpunan !!{sentence}s yang mempunyai interpretasi yang dimaksudkan,
entah tertutup secara deduktif ataupun tidak.

\begin{prop}
Bagi sembarang $\Struct{M}$, $\Theory{M}$ lengkap.
\end{prop}

\begin{proof}
Bagi sembarang !!{sentence}~$!A$, berlaku $\Sat{M}{!A}$ atau
$\Sat{M}{\lnot !A}$, sehingga $!A \in \Theory{M}$ atau
$\lnot !A \in \Theory{M}$.
\end{proof}

\begin{prop}\ollabel{prop:equiv}
  Jika $\Struct{N} \models !A$ bagi setiap $!A \in \Theory{M}$, maka
  $\Struct{M} \elemequiv \Struct{N}$.
\end{prop}

\begin{proof}
Karena $\Sat{N}{!A}$ bagi semua $!A \in \Theory{M}$, berlaku
$\Theory{M} \subseteq \Theory{N}$. Jika $\Sat{N}{!A}$, maka
$\Sat/{N}{\lnot !A}$. Karena menurut hipotesis struktur tersebut memenuhi
setiap kalimat dari teori yang pertama, hal ini menyiratkan
$\lnot !A \notin \Theory{M}$. Karena
$\Theory{M}$ lengkap, berlaku $!A \in \Theory{M}$. Jadi,
$\Theory{N} \subseteq \Theory{M}$, dan dengan demikian
$\Struct{M} \elemequiv \Struct{N}$.
\end{proof}

\begin{rem}\ollabel{remark:R}
  Tinjau $\Struct{R} = \langle\Real, <\rangle$, yakni !!{structure} yang
  domainnya merupakan himpunan bilangan real~$\Real$, dalam !!{language}
  yang hanya memuat satu !!{predicate} beraritas~2 yang ditafsirkan sebagai
  relasi~$<$ pada bilangan real. Jelas bahwa $\Struct{R}$
  !!{nonenumerable}; namun, karena $\Theory{R}$ jelas konsisten, menurut
  teorema L\"owenheim--Skolem teori tersebut mempunyai model yang
  !!{enumerable}, sebut saja $\Struct{S}$, dan menurut \olref{prop:equiv},
  $\Struct{R} \equiv \Struct{S}$. Selain itu, karena $\Struct{R}$ dan
  $\Struct{S}$ tidak isomorfik, hal ini menunjukkan bahwa kebalikan dari
  \olref[iso]{thm:isom} pada umumnya tidak berlaku.
\end{rem}

\end{document}
